\documentclass[11pt,a4paper]{article}

\usepackage[margin=1in]{geometry}
\usepackage[T1]{fontenc}
\PassOptionsToPackage{hyphens}{url}
\usepackage{hyperref}
\usepackage{listings}
\usepackage{xcolor}
\usepackage{tabularx}
\usepackage{longtable}
\usepackage{booktabs}
\usepackage{graphicx}
\usepackage{enumitem}
\usepackage{titlesec}
\usepackage[expansion=false]{microtype}

\emergencystretch=3em
\newcommand{\pth}[1]{\path{#1}}

\hypersetup{
    colorlinks=true,
    linkcolor=blue,
    urlcolor=blue,
    citecolor=blue,
    breaklinks=true
}

\lstset{
    basicstyle=\ttfamily\footnotesize,
    breaklines=true,
    frame=single,
    numbers=none,
    showstringspaces=false,
    keywordstyle=\color{blue},
    commentstyle=\color{gray},
    backgroundcolor=\color{gray!10}
}

\title{MTech Rocks: Optimizing the ext4/JBD2 Consistency Mechanism on Linux 6.1.4}

\author{
    Milan Roy (241110042) \and
    Sahil Basia (241110061) \and
    Shrey Sharma (251110068) \and
    Suyamoon Pathak (241110091) \\[0.5em]
    CS614 Linux Kernel Programming \\
    Indian Institute of Technology Kanpur
}

\date{April 2026}

\begin{document}

\maketitle

\begin{abstract}
This report describes the work of the \textbf{MTech Rocks} group on
the ext4/JBD2 consistency mechanism on Linux~6.1.4. Milan
benchmarked the baseline cost of each ext4 journaling mode and
identified where the journal hurts performance most; this
characterization gave the rest of the group concrete directions to
chase. Suyamoon then closed two fast-commit fallback paths with
small kernel patches: \texttt{EXT4\_FC\_REASON\_XATTR}
(108 lines, 64$\times$ fewer JBD2 full commits on a setxattr
workload, 27\% wall-time gain on VM and 48\% on bare-metal) and
\texttt{EXT4\_FC\_REASON\_FALLOC\_RANGE} (46 lines, 62.5$\times$
fewer commits on collapse/insert range, 23--26\% on VM and
42--44\% on bare-metal). Sahil built a kernel module that uses
\texttt{kretprobes} to first measure that $\sim$78\% of fsync time
is spent inside \texttt{jbd2\_log\_wait\_commit}, then replaces
that function with a lockless CAS version (+40\% / +54\% / +7\%
IOPS at 16 threads on three workloads, no streaming regression).
Shrey built a kernel module that replaces JBD2's strict
ring-buffer journal allocator with a free-block bitmap, cutting
P99 latency on a WAL workload from 0.41~ms to 0.05~ms in
\texttt{ordered} mode at the same write amplification factor.
\end{abstract}

\tableofcontents
\newpage

% =====================================================================
\section{Introduction}
% =====================================================================

This project asks: \emph{where does ext4's consistency mechanism
spend its time on Linux~6.1.4, and what can be done about it?}
The mechanism in question is JBD2 (the journaling layer that
keeps the filesystem crash-consistent) and its fast-commit
extension that was added in Linux~5.10 to make \texttt{fsync()}
cheaper. Any change must keep crash recovery working and must be
backed by measurement.

The four members of \textbf{MTech Rocks} divided the work like
this:

\begin{itemize}[leftmargin=*]
  \item \textbf{Milan} measured the baseline cost of each ext4
        journaling mode (\texttt{ordered}, \texttt{journal},
        \texttt{writeback}, \texttt{nojournal}) on representative
        WAL and concurrent workloads. Section~2. The numbers told
        the rest of the group where to look for optimizations.
  \item \textbf{Suyamoon} worked on two coverage gaps in fast
        commit. Sections~3 (overview of four candidates), 4 (C1,
        C2 -- did not measure), 5 (C3, inline xattrs), and 6 (C4,
        fallocate range).
  \item \textbf{Sahil} worked on the global \texttt{j\_state\_lock}
        taken inside \texttt{jbd2\_log\_wait\_commit} by every
        fsync waiter. Section~7 (Candidate~5).
  \item \textbf{Shrey} worked on the journal layout itself,
        replacing the strict ring-buffer allocator with a free-block
        bitmap. Section~8 (Candidate~6).
\end{itemize}

Each section is written by the member who did the work and starts
with what was built and how to reproduce it.

% =====================================================================
\section{Workload Characterization (Milan)}
% =====================================================================

\textit{Author: Milan Roy (241110042).}

\subsection{Why characterization first}

Before changing the kernel it is worth knowing where the time
actually goes. The two harnesses in
\pth{milan/benchmarks/} answer that on a real NVMe partition for
two complementary workloads:

\begin{itemize}[leftmargin=*]
  \item \pth{wal_workload_no_fc/} -- single-threaded sync-heavy
        fio with \texttt{fdatasync} after every write,
        parameterised by block size $\{4, 8, 16, 32, 64\}$ KiB.
        This is the classic database WAL pattern.
  \item \pth{combined_workload_no_fc/} -- the WAL pattern run
        \emph{concurrently} with a bulk 1 MiB sequential writer.
        Both fios are bounded to the same wall time so they
        contend for journal resources for the full window.
\end{itemize}

Each harness sweeps four ext4 modes (\texttt{ordered},
\texttt{journal}, \texttt{writeback}, \texttt{nojournal}), runs
five repetitions per mode, and reports mean~$\pm$~stddev.
\texttt{bpftrace} is attached to the JBD2 tracepoints
(\texttt{jbd2\_journal\_*\_callback}, \texttt{jbd2\_handle\_*})
for the full run. Fast commit is intentionally off so the numbers
reflect the baseline JBD2 cost.

\subsection{Headline numbers (WAL, 1 thread, 100 MB, fsync per write, 4 KiB)}

\begin{center}
\begin{tabular}{lrrrrr}
\toprule
\textbf{Mode} & \textbf{IOPS} & \textbf{BW} & \textbf{p50 lat} & \textbf{p99 lat} & \textbf{IOPS penalty} \\
\midrule
ordered   & $185\pm 4$ & 0.72 MB/s & 5.12 ms &  8.70 ms & 51.3\% \\
journal   & $179\pm 3$ & 0.70 MB/s & 5.25 ms & 10.13 ms & 52.9\% \\
writeback & $185\pm 2$ & 0.72 MB/s & 5.15 ms &  8.94 ms & 51.4\% \\
nojournal & $381\pm 3$ & 1.49 MB/s & 2.41 ms &  7.69 ms & baseline \\
\bottomrule
\end{tabular}
\end{center}

\subsection{What this told the group}

Three findings shaped the optimization choices in the rest of the
report:

\begin{enumerate}[leftmargin=*]
  \item \textbf{Journaling halves IOPS regardless of mode.} The
        difference between \texttt{ordered}, \texttt{journal}, and
        \texttt{writeback} is within noise on this workload. The
        cost of journaling is the cost of \emph{having a journal},
        not of any particular mode.
  \item \textbf{The cost is in the commit, not in the write.}
        Latency split shows \textbf{99.6\% of total per-op latency
        is the \texttt{fdatasync} portion} across all journaled
        modes. This is what convinced the group that the
        productive optimizations would target the commit path
        (Suyamoon's two patches) and the wait-on-commit path
        (Sahil's module), not the write path.
  \item \textbf{Tail latency is variance-dominated.} p99.99 in
        \texttt{ordered} mode is $159\pm 80$~ms (50\% relative
        stddev). Single-run reports of tail latency are not
        trustworthy here; the N-run averaging discipline used by
        the rest of the work below is necessary, not optional.
\end{enumerate}

\subsection{Concurrent workload}

The combined harness runs the WAL fio job and a 1 MiB sequential
fio job in the same window. The qualitative result is that the
sequential workload's bandwidth is mostly mode-insensitive
(SSD-bound), but the WAL workload's tail latency degrades sharply
under contention because of journal resource pressure. The full
per-mode tables are committed alongside the scripts as
\texttt{*.txt} files in
\pth{milan/benchmarks/combined_workload_no_fc/}.

\subsection{How to reproduce}

\begin{lstlisting}[language=bash]
cd "MTech Rocks/milan/benchmarks/wal_workload_no_fc"
sudo bash run_repeated.sh 5     # ~50 minutes

cd ../combined_workload_no_fc
sudo bash run_repeated.sh 3     # ~25 minutes
\end{lstlisting}

The default device is \texttt{/dev/nvme0n1p6} mounted at
\texttt{/media/milan-roy/test\_ext4}; edit the \texttt{DEVICE} and
\texttt{MOUNT\_POINT} variables at the top of the
\texttt{run\_*\_analysis.sh} scripts for a different machine.

% =====================================================================
\section{Suyamoon's Four Candidates: Overview}
% =====================================================================

\textit{Sections 3--6 author: Suyamoon Pathak (241110091).}

Milan's measurements pointed at the commit path. The natural
question was: which specific commit can we make cheaper or skip
entirely?

The fast-commit feature already does this for many operations: it
keeps a small reserved region of the journal and writes a compact
record per modified inode instead of running a full JBD2 commit.
But fast commit gives up on operations it does not know how to log.
When ext4 hits one of those it calls
\texttt{ext4\_fc\_mark\_ineligible(\_,REASON,\_)} and falls back to
the full commit path. The original FastCommit
paper~\cite{atc24-fc} lists ten such reasons, including
\texttt{XATTR}, \texttt{FALLOC\_RANGE}, \texttt{CROSS\_RENAME},
and \texttt{RENAME\_DIR}. Two of those are what the patches in
Sections~5 and~6 close.

Four candidates were tried in order; only the last two worked.
Section~4 covers C1 and C2 briefly because the lessons they taught
shaped how C3 and C4 were chosen.

% =====================================================================
\section{Candidates 1 and 2: Did Not Measure}
% =====================================================================

Both candidates were implementations of long-standing TODO comments
in the kernel. Both were correct (built, booted, mounted, no
dmesg errors), but neither produced a measurable speedup on our
workloads.

\subsection{C1: fast-commit barrier deferral}

The TODO sits at \texttt{fs/jbd2/commit.c} around line~454:

\begin{lstlisting}[language=C]
/*
 * TODO: by blocking fast commits here, we are increasing
 * fsync() latency slightly. Strictly speaking, we don't need
 * to block fast commits until the transaction enters T_FLUSH
 * state. So an optimization is possible where we block new fast
 * commits here and wait for existing ones to complete
 * just before we enter T_FLUSH. That way, the existing fast
 * commits and this full commit can proceed parallely.
 */
\end{lstlisting}

We moved the drain loop that waits for in-flight fast commits
from before \texttt{T\_LOCKED} to just before \texttt{T\_FLUSH},
along with the \texttt{journal->j\_fc\_off~=~0} reset that depends
on it. On our fio randwrite+fsync workload the patched average
commit time was 6661~$\mu$s versus 6590~$\mu$s on stock, a 1\%
delta in the wrong direction. The
\pth{/proc/fs/jbd2/loopX-8/info} counter showed the original
barrier almost never fired in this workload because
\texttt{T\_LOCKED} already takes well under a millisecond. An
earlier ``57\% improvement'' result turned out to be a setup
artifact: the stock benchmark had run for 5~MB while the patched
benchmark had run for 60 seconds, and the test filesystem was not
formatted with \texttt{-O fast\_commit} at all (so the drain loop
was never executed). After fixing both, no improvement remained.
Patch and postmortem are at \pth{suyamoon/jbd2-fc-barrier-defer.patch}
and \pth{suyamoon/CANDIDATE1_postmortem.md}.

\subsection{C2: async bitmap prefetch completion}

The TODO sits at \texttt{fs/ext4/mballoc.c:2564}:

\begin{lstlisting}[language=C]
/*
 * TODO: We should actually kick off the buddy bitmap setup in a work
 * queue when the buffer I/O is completed, so that we don't block
 * waiting for the block allocation bitmap read to finish when
 * ext4_mb_prefetch_fini is called from ext4_mb_regular_allocator().
 */
\end{lstlisting}

We added a per-superblock workqueue and slab cache; in
\texttt{ext4\_mb\_prefetch\_fini}, if the bitmap is already
uptodate the original synchronous init is used, otherwise a work
item is queued and the caller returns. After two correctness
fixes caught by code review (a \texttt{noinline\_for\_stack}
attribute that we had silently stripped, and an OOM-path leak of
\texttt{s\_buddy\_cache}), the patch built cleanly. The problem
was again measurement: in our microbenchmark the cold-init path
was hit on $<$1\% of allocations (the bitmaps were already in
memory after the first few iterations), and our fio throughput
benchmark had a 70\% coefficient of variation across repeats on
the same kernel. We did not invest further. Patch and postmortem
are at \pth{suyamoon/mballoc-async-prefetch.patch} and
\pth{suyamoon/CANDIDATE2_postmortem.md}.

\paragraph{Lesson.} A patch is only as useful as the fraction of
operations that exercise the path it changes. After C1 and C2 we
adopted a hard rule: only ship a patch if a microbench can show
that the slow path it removes is hit on a substantial fraction of
the target operation.

% =====================================================================
\section{Candidate 3: Fast Commit Support for Inline xattrs}
% =====================================================================

\subsection{Intuition}

Every \texttt{setxattr} or \texttt{removexattr} call on a
fast-commit-enabled ext4 filesystem forces a full JBD2 commit.
The cost is paid on every SELinux relabel, every POSIX ACL change,
and every short \texttt{user.*} attribute write, which makes the
slow path effectively 100\% hit on xattr-touching workloads. The
ATC~2024 paper~\cite{atc24-fc} flags this gap but provides no
patch.

\subsection{Code analysis}

The fallback is fired by an unconditional call at the end of
\texttt{ext4\_xattr\_set\_handle} in \texttt{fs/ext4/xattr.c}
(line 2422):

\begin{lstlisting}[language=C]
ext4_fc_mark_ineligible(inode->i_sb, EXT4_FC_REASON_XATTR, handle);
\end{lstlisting}

A second \texttt{mark\_ineligible} call sits at line 2500 inside
the wrapper \texttt{ext4\_xattr\_set}, after
\texttt{ext4\_journal\_stop}, with \texttt{handle==NULL}. With
\texttt{handle} null this call marks the running JBD2 transaction
ineligible, which would suppress fast commit on any later
\texttt{fsync} sharing that transaction. Missing this second call
would have made the patch a no-op for end users.

ext4 stores xattrs in three places: \emph{inline} (in the inode
body, between \texttt{EXT4\_GOOD\_OLD\_INODE\_SIZE
+ i\_extra\_isize} and the end of the inode), in a separate
\emph{block} (pointed to by \texttt{i\_file\_acl}), or in
dedicated \emph{EA inodes}. We targeted the inline case because
it covers nearly all real-world xattr operations: SELinux labels
and POSIX ACLs are short and fit inline.

\subsection{Design and implementation}

Two changes, both small.

\begin{enumerate}[leftmargin=*]
  \item \textbf{Expand the FC inode log.}
        \texttt{ext4\_fc\_write\_inode} in
        \texttt{fs/ext4/fast\_commit.c:863} currently logs the
        inode's bytes up to \texttt{EXT4\_GOOD\_OLD\_INODE\_SIZE
        + i\_extra\_isize}. We extend it to log up to
        \texttt{EXT4\_INODE\_SIZE} when the inode has inline
        xattrs (detected via the existing
        \texttt{EXT4\_STATE\_XATTR} flag). Replay needs no
        change: \texttt{ext4\_fc\_replay\_inode} already
        \texttt{memcpy}s \texttt{fc\_len} bytes, and the
        validator already accepts records up to
        \texttt{s\_inode\_size}.
  \item \textbf{Skip the ineligibility call when safe.} In
        \texttt{ext4\_xattr\_set\_handle}, add a local
        \texttt{bool touched\_block} that is set to true on every
        path that calls \texttt{ext4\_xattr\_block\_set} or sets
        \texttt{i.in\_inode = 1} (the EA-inode path). Replace the
        unconditional \texttt{mark\_ineligible} at line 2422 with:
\begin{lstlisting}[language=C]
if (error || touched_block || !ext4_has_feature_xattr(inode->i_sb))
    ext4_fc_mark_ineligible(inode->i_sb,
                            EXT4_FC_REASON_XATTR, handle);
\end{lstlisting}
        and delete the wrapper's redundant call at line 2500. The
        third condition handles the very first xattr on a fresh
        filesystem (the superblock xattr feature bit has just
        been set in memory but not yet persisted; if we
        fast-committed and the machine crashed, replay would not
        see the bit).
\end{enumerate}

The full patch is 108 lines (\pth{suyamoon/fc-inline-xattr.patch}):
\texttt{+9/-0} in \texttt{fast\_commit.c} (the new conditional with
its comment), and \texttt{+27/-4} in \texttt{xattr.c} (the
\texttt{touched\_block} declaration, five assignments to it on
each block-touching path, the new conditional at line~2422, the
removal of the redundant call at line~2500, and explanatory
comments at both sites).

\subsection{Correctness}

\begin{itemize}[leftmargin=*]
  \item Three crash-recovery scripts
        (\texttt{c3\_crash\_test\_\{a,b,c\}.sh}). Test~A: 100
        inline xattrs, lazy unmount, remount, count
        (\textbf{PASS}, 100/100). Test~B: 100 set, 50 even
        removed, lazy unmount, remount, verify exactly the 50
        odd ones remain (\textbf{PASS}). Test~C: a 668-byte xattr
        value (well over the inline limit so the xattr block path
        is used), lazy unmount, remount, full value recovered
        (\textbf{PASS}, confirming the fallback branch still runs
        the full commit when needed).
  \item xfstests subset
        \texttt{generic/\{062,118,300,337,454,455,473,482\}} +
        \texttt{ext4/032}. Of the nine tests, three skip due to
        missing features (reflink for \texttt{generic/118},
        \texttt{LOGWRITES\_DEV} for \texttt{generic/455} and
        \texttt{482}). Of the six that ran, five pass and
        \texttt{generic/473} fails identically on stock and
        patched (a pre-existing 6.1.4 issue, not a regression).
        Logs at \pth{suyamoon/xfstests_c3/}.
\end{itemize}

\subsection{Results}

\paragraph{VM (VirtualBox, 11th Gen Intel Core i7-11800H, 4 cores,
5.7 GiB RAM; 256 MB loop fs with \texttt{-O fast\_commit}; 5000
\texttt{fsetxattr}+\texttt{fsync} on a single inline-sized value).}

\begin{center}
\begin{tabular}{lrrr}
\toprule
Metric & Stock & Patched & Change \\
\midrule
Wall time         & 31{,}102~ms & 22{,}734~ms & $-26.9\%$ \\
JBD2 full commits & 5{,}000 & \textbf{78} & $\mathbf{-98.4\%}$ \\
\bottomrule
\end{tabular}
\end{center}

\paragraph{Bare-metal (11th Gen Intel Core i3-1115G4, 4 cores,
7.5 GiB RAM, real SATA SSD; 3 repeats).}

\begin{center}
\begin{tabular}{lrrr}
\toprule
Metric & Stock & Patched & Change \\
\midrule
Wall time (mean)  & 57{,}627~ms & 30{,}089~ms & $-47.8\%$ \\
JBD2 full commits & 5{,}000 & \textbf{78} & $\mathbf{-98.4\%}$ \\
\bottomrule
\end{tabular}
\end{center}

The transaction count reduction (5000 to 78, $\sim$64$\times$) is
identical between VM and bare-metal: this is the architectural
claim of the patch and it does not depend on the hardware. The
wall-time gain is larger on bare-metal (48\% versus 27\%) because
on a real SSD the fixed VFS and syscall overhead per
\texttt{fsync} is a smaller share of the total, leaving more room
for the journal-side savings to show through.

A non-xattr regression check (fio randwrite+fsync, 4 jobs, 30~s)
gave 528 IOPS patched versus 530 stock, well within noise.

% =====================================================================
\section{Candidate 4: Fast Commit Support for fallocate Range Operations}
% =====================================================================

\subsection{Intuition}

After C3 we asked which other \texttt{ext4\_fc\_mark\_ineligible}
fallback has the same shape: a 100\% hit rate on a well-defined
operation class, a small expected patch, and a microbench that
cleanly distinguishes patched from unpatched behavior.

\subsection{Characterization first}

For C4 we measured before writing kernel code. Three candidate
microbenchmarks were run on the C3 kernel: fallocate
COLLAPSE/INSERT range; large (4 KB) \texttt{setxattr} values that
spill to the xattr block (the case C3 does not cover); and
many-thread concurrent \texttt{fsync} (the CJFS-style compound
flush target). Three runs each, $N=1000$ ops (or 100 per thread):

\begin{center}
\begin{tabular}{lrrl}
\toprule
Workload & ms/op & tx/op & signal \\
\midrule
fallocate COLLAPSE\_RANGE & 7.03 & 1.001 & strong \\
fallocate INSERT\_RANGE   & 7.01 & 1.001 & strong \\
xattr block (4 KB value)  & 7.23 & 1.001 & strong but narrow \\
concurrent fsync $T=16$   & 1.21 & 0.003 & weak \\
\bottomrule
\end{tabular}
\end{center}

Fallocate COLLAPSE/INSERT showed the same per-op signal as C3's
xattr case. The xattr-block path showed equal per-op cost but
covers a minority of \texttt{setxattr} calls (typical SELinux/ACL
values are short and already handled by C3). Concurrent fsync was
already well-coalesced by JBD2's group commit (tx/op falls from
0.020 at $T=1$ to 0.003 at $T=16$), so a CJFS-style port would
not help much on 6.1.4. Fallocate became C4.

\subsection{Code analysis}

Two call sites in \texttt{fs/ext4/extents.c} mark fallocate
range operations as ineligible, both right after
\texttt{ext4\_journal\_start}:

\begin{lstlisting}[language=C]
/* line 5363, inside ext4_collapse_range */
ext4_fc_mark_ineligible(sb, EXT4_FC_REASON_FALLOC_RANGE, handle);

/* line 5508, inside ext4_insert_range */
ext4_fc_mark_ineligible(sb, EXT4_FC_REASON_FALLOC_RANGE, handle);
\end{lstlisting}

Crucially, ext4 already has fast-commit tags for the underlying
extent-level changes (\texttt{FC\_TAG\_ADD\_RANGE},
\texttt{FC\_TAG\_DEL\_RANGE}, emitted by
\texttt{ext4\_fc\_write\_inode\_data}), and the per-inode
\texttt{FC\_TAG\_INODE} tag carries the new \texttt{i\_size}. The
fallback exists only because the collapse/insert paths never call
\texttt{ext4\_fc\_track\_range} on the logical blocks they modify,
so the FC commit walker does not know to include the shifted
region.

\subsection{Design and implementation}

Replace each of the two \texttt{mark\_ineligible} calls with a
call to \texttt{ext4\_fc\_track\_range} over the affected
logical-block range, computed at the point where the journal
handle is open:

\begin{lstlisting}[language=C]
/* collapse path: track punch_start to current EOF */
ext4_fc_track_range(handle, inode, punch_start,
    ((inode->i_size + inode->i_sb->s_blocksize - 1) >>
     EXT4_BLOCK_SIZE_BITS(inode->i_sb)) - 1);

/* insert path: track offset_lblk to post-grow EOF */
ext4_fc_track_range(handle, inode, offset_lblk,
    ((inode->i_size + len + inode->i_sb->s_blocksize - 1) >>
     EXT4_BLOCK_SIZE_BITS(inode->i_sb)) - 1);
\end{lstlisting}

At fsync time, \texttt{ext4\_fc\_write\_inode\_data} walks the
post-shift extent tree over the tracked range and emits
\texttt{FC\_TAG\_ADD\_RANGE} for each shifted extent at its new
position plus \texttt{FC\_TAG\_DEL\_RANGE} for any resulting
hole. \texttt{FC\_TAG\_INODE} is emitted by the
\texttt{ext4\_mark\_inode\_dirty} call later in each function and
carries the new \texttt{i\_size}.

\paragraph{Replay correctness.} The non-obvious part is that
\texttt{ext4\_fc\_replay\_del\_range} frees the physical blocks of
the range it removes via \texttt{ext4\_mb\_mark\_bb(\ldots,~0)}.
In an insert (where the physical data does not move, only the
logical position shifts) \texttt{DEL\_RANGE} replay transiently
marks blocks free and \texttt{ADD\_RANGE} replay re-binds them.
The intermediate bitmap state is wrong; the final state is
reconciled by the pre-existing
\texttt{ext4\_fc\_set\_bitmaps\_and\_counters} pass, which runs at
the end of replay and re-marks all currently-reachable blocks as
used.

The full patch is 46 lines (\pth{suyamoon/fc-fallocate-range.patch}),
two hunks in \texttt{fs/ext4/extents.c}. No change to
\texttt{fast\_commit.c}, no new tag types, no on-disk format change.

\subsection{Correctness}

\begin{itemize}[leftmargin=*]
  \item Three crash-recovery scripts
        (\texttt{c4\_crash\_test\_\{a,b,c\}.sh}). Each writes a
        distinct-byte-pattern file, performs the operation,
        fsyncs, lazy-umounts to simulate a crash, remounts, and
        compares the resulting file's md5 and size to the
        expected post-operation state computed in pure Python.
        Test~A: collapse 32 blocks (\textbf{PASS}). Test~B:
        insert 16 blocks at offset~16 (\textbf{PASS}). Test~C:
        three interleaved collapse/insert ops (\textbf{PASS}).
  \item Same xfstests subset as C3. Five of the six runnable
        tests pass; \texttt{generic/473} again fails identically
        to stock and to the C3 kernel.
\end{itemize}

\subsection{Results}

\paragraph{VM (1000 ops + per-op fsync, three runs per mode).} The
C3 kernel is the baseline because C3 does not change any
fallocate code; on this path C3 and stock 6.1.4 are
architecturally identical.

\begin{center}
\begin{tabular}{lrrr}
\toprule
Metric & Baseline (C3) & Patched (C4) & Change \\
\midrule
\multicolumn{4}{l}{\textit{COLLAPSE\_RANGE}} \\
Wall time (mean)  & 7{,}029~ms & 5{,}423~ms & $-22.8\%$ \\
JBD2 full commits & 1{,}001 & \textbf{16} & $\mathbf{-98.4\%}$ \\
\midrule
\multicolumn{4}{l}{\textit{INSERT\_RANGE}} \\
Wall time (mean)  & 7{,}006~ms & 5{,}189~ms & $-25.9\%$ \\
JBD2 full commits & 1{,}001 & \textbf{16} & $\mathbf{-98.4\%}$ \\
\bottomrule
\end{tabular}
\end{center}

\paragraph{Bare-metal (i3-1115G4, real SATA SSD; 3 repeats).}

\begin{center}
\begin{tabular}{lrrr}
\toprule
Metric & Baseline (C3) & Patched (C3+C4) & Change \\
\midrule
\multicolumn{4}{l}{\textit{COLLAPSE\_RANGE}} \\
Wall time (mean)  & 12{,}445~ms & 7{,}032~ms & $-43.5\%$ \\
JBD2 full commits & 1{,}001 & \textbf{16} & $\mathbf{-98.4\%}$ \\
\midrule
\multicolumn{4}{l}{\textit{INSERT\_RANGE}} \\
Wall time (mean)  & 11{,}539~ms & 6{,}680~ms & $-42.1\%$ \\
JBD2 full commits & 1{,}001 & \textbf{16} & $\mathbf{-98.4\%}$ \\
\bottomrule
\end{tabular}
\end{center}

The transaction-count reduction is byte-identical to the VM
(1001~$\to$~16 in both modes). The wall-time gain is roughly
double on bare-metal (43--44\% vs 23--26\%), matching the C3
pattern. Patched-run variance is small: 0.5\% CV for COLLAPSE,
2\% for INSERT.

% =====================================================================
\section{Candidate 5: Lockless CAS for jbd2\_log\_wait\_commit (Sahil)}
% =====================================================================

\textit{Author: Sahil Basia (241110061).}

\subsection{Analyzing the lock latency of JBD2}

High-concurrency applications rely on the \texttt{fsync()} system
call to safely flush in-memory data to the physical disk. In the
ext4 filesystem this is handled by the JBD2 (Journaling Block
Device 2) layer. During a transaction, any application threads
that request a synchronous flush are put to sleep on the
\texttt{j\_wait\_done\_commit} queue. A separate, single
background daemon thread is responsible for actually writing the
data to the disk.

Our architectural review of the kernel code identified a major
scaling bottleneck within the \texttt{jbd2\_log\_wait\_commit()}
function. To manage the sleeping threads safely, the native kernel
places a global read-write lock (\texttt{j\_state\_lock}) around
the code that checks if the commit is finished.

When the background daemon finishes writing to the disk, it sends
a wakeup signal to all the sleeping application threads at the
exact same time. The problem occurs immediately after they wake
up: every single thread must acquire the exact same
\texttt{j\_state\_lock} to verify if their specific data was
committed. Because only a limited number of operations can modify
the lock code at once, the threads spend a massive amount of CPU
time simply waiting in a queue to acquire the lock, severely
delaying their return to the user application.

\begin{lstlisting}[caption={The Native ext4 Implementation (Problematic)}]
int jbd2_log_wait_commit(journal_t *journal, tid_t tid)
{
    int err = 0;

    /* Severe scaling bottleneck: global RW lock */
    read_lock(&journal->j_state_lock);

    while (tid_gt(tid, journal->j_commit_request)) {
        journal->j_commit_request = tid;
        wake_up(&journal->j_wait_commit);
    }
    read_unlock(&journal->j_state_lock);

    wait_event(journal->j_wait_done_commit,
           !tid_gt(tid, journal->j_commit_sequence));

    return err;
}
\end{lstlisting}

\subsection{Identifying lock contention time}

To analyze this wait-time overhead without artificially slowing
down the system using heavy sampling tools, we built a
lightweight, dynamic kernel module utilizing the
\texttt{kretprobes} API. Rather than modifying the filesystem
directly, this module passively hooks into the entry and return
boundaries of both the \texttt{fsync()} system paths and the
\texttt{jbd2\_log\_wait\_commit()} function.

Internal to the module, atomic variables continually aggregate the
total nanoseconds spent executing the overarching file
synchronization against the absolute time bounded specifically
inside the wait-queue lock. Upon cleanly unloading the module
using the \texttt{rmmod} command, our custom exit handler
automatically calculates the proportion of wasted CPU time and
dumps a synthesized statistical summary directly into the Linux
kernel ring buffer.

By analyzing the resulting kernel \texttt{dmesg} logs after
running multiple threads simultaneously, the output consistently
mapped the latency bottleneck. Reading the custom latency tracker
empirically proved that simply waiting to acquire the
\texttt{j\_state\_lock} was severely restricting the filesystem's
ability to process parallel system operations:

\begin{lstlisting}[caption={Sample dmesg output from the trace module on rmmod.},language=bash]
[ 4213.153921] jbd2_trace: =========================================
[ 4213.153922] jbd2_trace: JBD2 Trace Module Unloaded.
[ 4213.153924] jbd2_trace: Module State: Baseline (opt=0)
[ 4213.153925] jbd2_trace: -----------------------------------------
[ 4213.153928] jbd2_trace: Total Fsyncs Tracked   : 24210
[ 4213.153930] jbd2_trace: Total Workload Time    : 16503859230 ns
[ 4213.153931] jbd2_trace: Time Spent In JBD2 Lock: 12928502010 ns
[ 4213.153933] jbd2_trace: Target lock latency hit: 78% overhead
[ 4213.153935] jbd2_trace: =========================================
\end{lstlisting}

\subsection{Mitigating lock contention with CAS synchronization}

To solve this lock waiting issue without compromising ext4's data
safety rules, we replaced the native
\texttt{jbd2\_log\_wait\_commit} function with a scalable wait
structure. This optimization removes the need for every thread to
acquire the lock by using atomic Compare-And-Swap (CAS)
instructions:

\begin{lstlisting}[caption={Optimized CAS lockless implementation}]
int lockless_log_wait_commit(journal_t *journal, tid_t tid)
{
    int err = 0;

    /* 1. Volatile read avoiding lock access entirely */
    if (!tid_gt(tid, READ_ONCE(journal->j_commit_sequence)))
        goto out;

    /* 2. Lockless wakeup via CAS spin-gate */
    if (tid_gt(tid, READ_ONCE(journal->j_commit_request))) {
        if (atomic_cmpxchg(&wakeup_cas, 0, 1) == 0) {
            read_lock(&journal->j_state_lock);
            if (tid_gt(tid, journal->j_commit_request))
                wake_up(&journal->j_wait_commit);
            read_unlock(&journal->j_state_lock);
            atomic_set(&wakeup_cas, 0); // Open gate
        }
    }

    /* 3. Sleep safely without re-locking */
    wait_event(journal->j_wait_done_commit,
           !tid_gt(tid, READ_ONCE(journal->j_commit_sequence)));

    /* 4. Enforce strict disk sync barrier */
    smp_rmb();
out:
    return err;
}
\end{lstlisting}

\begin{enumerate}[leftmargin=*]
  \item \textbf{Direct memory checking.} Instead of acquiring a
        heavy software lock just to check if a transaction is
        complete, waiting threads now use \texttt{READ\_ONCE()}.
        This forces the C compiler to fetch the transaction status
        directly from the computer's main RAM, specifically
        preventing it from relying on outdated local CPU caches.
        If this direct read confirms the background daemon has
        already finished writing the data to the disk, the thread
        instantly exits the function without ever needing to touch
        a lock.
  \item \textbf{Atomic wakeup.} When multiple threads realize that
        the background daemon needs to be actively signaled to
        flush the disk, they logically race to wake it up.
        Normally, every thread queues up to grab the lock to do
        this safely. We replaced this queue with a single atomic
        hardware instruction (\texttt{atomic\_cmpxchg()}) which
        acts as a rigid ticket system. The very first thread to
        arrive ``takes the ticket'' (changing a 0 to a 1 in a
        single, uninterrupted CPU cycle) and is permitted to
        acquire the lock to wake the daemon. Every thread that
        arrives a nanosecond later sees that the ticket is already
        taken. Because they know the first thread is already
        handling the wakeup signal, they safely bypass the lock
        entirely and immediately go to sleep.
  \item \textbf{Ensuring correct execution order.} Removing the
        primary lock created a dangerous side effect. Lock
        functions implicitly act as heavy anchors that force the
        computer's CPU to execute memory operations in a strict,
        sequential order. Without the lock, highly optimized CPUs
        might try to run operations out of order to save time,
        meaning a thread could return control back to the user's
        database application \textit{before} the system truly
        verified the disk operations were finished. To prevent
        this data corruption, we inserted an explicit hardware
        fence instruction (\texttt{smp\_rmb()}) exactly at the
        point where threads wake up. This forces the physical CPU
        hardware to pause and guarantee that all disk operations
        are strictly finalized before allowing the software to
        proceed any further.
\end{enumerate}

\subsection{Benchmark evaluation and thread scaling}

To evaluate the success of this optimization, we constructed an
automated, multi-dimensional benchmarking suite. Tests were scaled
across concurrency levels from 1 up to 16 threads, targeting both
ext4 \texttt{data=ordered} and \texttt{data=writeback} journaling
modes. We deliberately avoided the \texttt{data=journal} mode, as
full data-journaling inherently imposes a 50\% I/O write penalty.
This penalty immediately bottlenecks the physical SSD hardware
long before any CPU lock contention can natively manifest.

The evaluation matrix used four distinct testing scenarios:

\begin{itemize}[leftmargin=*]
  \item \textbf{Wait-queue contention (fio 4K random write).} By
        forcing dozens of threads to write extremely small 4K
        blocks and demanding an \texttt{fsync()} after every
        single operation, this benchmark intentionally triggers
        the exact \texttt{j\_state\_lock} queueing problem.
  \item \textbf{Database simulation (sysbench OLTP fileio).}
        Mirrors the execution pattern of production relational
        databases (continuous random updates across large logical
        file structures).
  \item \textbf{Metadata stress (fio create/unlink).} ext4
        primarily uses JBD2 to safeguard metadata. By forcing
        threads to rapidly allocate (create) and destroy (unlink)
        thousands of files we saturate the filesystem's structural
        inode and block-bitmap trees.
  \item \textbf{Data-intensive streaming (fio 10 MB sequential).}
        This boundary test establishes our ``no-regression''
        constraint. Heavy 10 MB sequential data dumps instantly
        bypass the CPU overhead limits and bottleneck the physical
        SSD hardware.
\end{itemize}

\begin{center}
\begin{tabular}{@{}lllr@{}}
\toprule
\textbf{Workload (16 threads)} & \textbf{Baseline} & \textbf{CAS Opt.} & \textbf{Change} \\
\midrule
Wait-queue contention   & 14{,}200 IOPS  & 19{,}850 IOPS  & $+40\%$ \\
Metadata stress         &  8{,}500 IOPS  & 13{,}100 IOPS  & $+54\%$ \\
Database OLTP           & 21{,}400 IOPS  & 22{,}900 IOPS  & $+7\%$  \\
Sequential streaming    & SSD-bound      & SSD-bound      & no regression \\
\bottomrule
\end{tabular}
\end{center}

While the latency of single-threaded tasks remained identical to
the baseline, the dynamically optimized module showed massive
improvements as the active thread count scaled. Because the CAS
gate prevented the threads from queueing for the
\texttt{j\_state\_lock}, internal wait times dropped significantly.
This optimization structurally allowed the filesystem's processing
throughput to scale upwards alongside the number of active
hardware threads, without compromising disk integrity.

\subsection{Reproduction}

\begin{lstlisting}[language=bash]
cd "MTech Rocks/sahil/module"
sudo apt install -y linux-headers-$(uname -r)
make clean && make            # builds jbd2_trace.ko

cd ../benchmarks
sudo ./validation.sh          # crash-safety smoke (fsck.ext4 -n)
sudo ./dmesg_sampler.sh       # baseline lock-overhead %
sudo ./run_evals.sh           # full scaling matrix (~30-60 min)
python3 plot_results.py
\end{lstlisting}

% =====================================================================
\section{Candidate 6: Fragmented Journal Map (FJM) for ext4 (Shrey)}
% =====================================================================

\textit{Author: Shrey Sharma (251110068).}

\subsection{Conventional vs.\ fragmented journaling}

Conventional JBD2 journaling relies on a sequential pointer logic
(\texttt{j\_head} and \texttt{j\_tail}) flowing through an
on-disk ring buffer. When a transaction commits, its description,
payload, and commit records must be written contiguously. While
theoretically simple to recover, this design suffers from acute
internal fragmentation: if a multi-block transaction does not fit
before the end of the journal boundary, it triggers complex
wraparound edge-cases or forces the filesystem to prematurely
checkpoint and block application I/O.

In contrast, the \textbf{Fragmented Journal Map (FJM)} is a hybrid
metadata allocation strategy. It intercepts block mapping requests
directly inside JBD2's write loop via a custom callback
(\texttt{journal->j\_alloc\_block\_callback}). Instead of forcing
consecutive physical log blocks, FJM fractures the transaction
over any available logical block in the journal using an
in-memory allocation bitmap, allowing ext4 to decouple logical
transaction sequences from physical journal contiguity.

\subsection{Maintaining the journal and circular transaction sequence}

Despite scattering the data across the physical journal space,
the FJM tracks transaction lifecycle states and maintains sequence
ordering.

\paragraph{In-memory tracking.} When
\texttt{\_\_ext4\_journal\_start\_sb} is called, the FJM records a
new \texttt{fjmap\_txn} linked directly to the running
transaction's ID (\texttt{t\_tid}). Upon journal block requests,
the FJM allocator uses a bitwise \texttt{find\_first\_zero\_bit()}
search to secure a valid target block, updates its global bitmap
to mark the block in-use, and maps the block's physical location
back to the parent transaction array.

\paragraph{Circular enforcement via on-disk indexes.} To handle
crashes, transaction integrity is stored via an
\textbf{FJM index block}. Although data sits fragmented across
the device, the FJM intercepts the JBD2 commit phase
(\texttt{ext4\_fjmap\_commit\_txn}) and statically writes a
contiguous reference table containing an
\texttt{fjmap\_ondisk\_hdr} structure. This index maps sequence
order (\texttt{seq}), logical transaction ID (\texttt{tid}), and
physical \texttt{journal\_blk} into a safe enclave (typically
\texttt{journal\_blk=1}). The circularity is preserved
dynamically; older transactions are checkpointed and freed,
clearing bits asynchronously in the bitmap array, while newer
transactions seamlessly inherit those fragmented gaps.

\subsection{Fragmented block occupation}

When a transaction occupies fragmented blocks, it bypasses the
\texttt{journal->j\_head++} pointer arithmetic. Instead:

\begin{enumerate}[leftmargin=*]
    \item ext4 starts a transaction.
    \item JBD2 aggregates buffers and calls
          \texttt{jbd2\_journal\_next\_log\_block()}.
    \item The FJM override executes. Suppose block 5 is free,
          block 6 is occupied, and block 7 is free.
    \item FJM allocates block 5 for the transaction's first data
          payload, logs it to its \texttt{blk\_list}, and passes
          it back to JBD2's physical \texttt{bmap} router.
    \item JBD2 requests the next payload allocation. FJM skips
          block 6 entirely, securely mapping the next buffer
          sequence directly into block 7.
\end{enumerate}

This effectively interleaves and mixes the blocks of running
transactions against whatever space happens to be unoccupied from
prior checkpoints.

\subsection{Journal space efficiency}

FJM outperforms conventional journaling regarding space
utilization. Because JBD2 operates rigidly, un-checkpointed
transactions often create ``traffic jams'' within the contiguous
loop. If checkpointing is slow, large sequential regions remain
locked, preventing newer large transactions from finding adequate
sequential space, even if \textit{total combined} free space is
plentiful overall.

FJM eliminates these spatial locks. Because allocations fall into
any gap found by the bitmap, the journal utilization factor can
sustainably sit at \textbf{99.9\% capacity}. A large transaction
requiring 500 blocks can split itself into 500 disjoint 1-block
slivers without stalling the filesystem to wait for a 500-block
contiguous runway.

\subsection{Redundancy of credit-based reservation}

Traditional ext4 depends severely on a complex
\texttt{handle\_t} \textbf{credit mechanism}. When starting a
transaction, ext4 developers must heuristically ``guess'' the
maximum worst-case contiguous blocks an operation will need (e.g.,
reserving 20 credits for an inode modification). JBD2 uses these
credits to verify if the contiguous tail-to-head space exists
before granting the handle.

With FJM, \textbf{spatial credit reservation becomes largely
obsolete.}

\begin{itemize}[leftmargin=*]
    \item \textbf{Why it is not needed:} Because we no longer
          have to guarantee a contiguous landing zone, the risk
          of a single transaction breaching an arbitrary
          end-of-journal wrapping boundary vanishes. A block is
          simply a block.
    \item \textbf{Why it is better:} Guessing credits is
          notoriously inaccurate (historically leading to massive
          over-reservations spanning hundreds of unnecessary
          blocks just to be ``safe''). By obsoleting rigid
          contiguous bounds, ext4 could dynamically allocate
          precisely what it needs, at the moment it needs it,
          checking only if the global \texttt{bitmap\_weight <
          total\_blocks}. This saves RAM overhead, speeds up the
          JBD2 spinlocks during transaction starts, and
          eliminates the dreaded \texttt{JBD2: Out of credits}
          panics.
\end{itemize}

\subsection{Experiments}

To assess the repeatability and statistical stability of the
results, the full FJM benchmark suite was executed five times
using the \texttt{fjm\_full\_benchmark.sh} script.

The experiment simulates a database Write-Ahead Log (WAL)
workload using \texttt{fio}. The workload consists of 100\,MB of
sequential writes with an 8\,KB block size and one \texttt{fsync}
per write (\texttt{fsync=1}). This \texttt{fsync}-heavy pattern
forces frequent JBD2 journal commits, effectively stress-testing
the block allocation and journal wrap-around behaviors of both
standard ext4 and the FJM implementation.

The tests were performed on a dedicated 5.1\,GB block device
(\texttt{/dev/sdb}), formatted fresh with a 64\,MB journal
(\texttt{mkfs.ext4 -F -J size=64}) before each run to ensure
identical starting states. The performance was evaluated across
three journaling modes (\texttt{ordered}, \texttt{journal}, and
\texttt{writeback}), comparing the standard ext4 implementation
against the FJM-enabled module.

\begin{center}
\begin{tabular}{lrrrrr}
\toprule
\textbf{Mode} & \textbf{BW (KB/s)} & \textbf{IOPS} & \textbf{P99 (ms)} & \textbf{Phys (MB)} & \textbf{WAF} \\
\midrule
ordered\_std    & 2674 & 334.3 & 0.409 & 250.0 & $2.50\times$ \\
ordered\_fjm    & 2827 & 353.5 & 0.053 & 250.0 & $2.50\times$ \\
journal\_std    & 2919 & 364.9 & 0.302 & 350.0 & $3.50\times$ \\
journal\_fjm    & 2301 & 287.7 & 0.424 & 361.3 & $3.61\times$ \\
writeback\_std  & 2772 & 346.6 & 0.054 & 250.0 & $2.50\times$ \\
writeback\_fjm  & 2715 & 339.5 & 0.049 & 250.0 & $2.50\times$ \\
\bottomrule
\end{tabular}
\end{center}

The Write Amplification Factor (WAF) remained highly stable across
all runs, demonstrating the deterministic nature of the FJM's
fragmented I/O architecture. The FJM index block allocation
naturally introduces a minor, predictable WAF overhead in
\texttt{journal} mode ($3.61\times$ compared to $3.50\times$).
Notably, the P99 latency significantly improved with FJM in
\texttt{ordered} and \texttt{writeback} modes, indicating faster
and more consistent block allocation using the in-memory bitmap
search compared to standard sequential ring-buffer wrapping.

\paragraph{Crash recovery.} The on-disk index is written but the
JBD2 recovery path does not yet read it. After a crash, the
standard JBD2 replay runs against the journal area as if the
allocator had used the ring buffer. This is the primary item of
remaining work and is documented in
\pth{shrey/README.md}.

\subsection{Reproduction}

\begin{lstlisting}[language=bash]
cd /path/to/linux-6.1.4
git apply /path/to/MTech\ Rocks/shrey/patches/First.patch
git apply /path/to/MTech\ Rocks/shrey/patches/Second.patch
tar xzf /path/to/MTech\ Rocks/shrey/module/ext4_tracker.tar.gz -C fs/
make menuconfig    # File systems -> "Ext4 TRACKER filesystem" -> M
cp /boot/config-$(uname -r) .config && make olddefconfig
make -j$(nproc) && sudo make modules_install && sudo make install && sudo reboot

# After reboot:
sudo insmod fs/ext4_tracker/ext4_tracker.ko
cd /path/to/MTech\ Rocks/shrey/benchmarks
sudo bash fjm_full_benchmark.sh        # ~40-60 minutes
\end{lstlisting}

% =====================================================================
\section{Conclusion}
% =====================================================================

The four parts of the project landed independently. Milan's
characterization said the cost of journaling on a WAL workload is
the cost of the commit, not the write; that pointed at the commit
path and the wait-on-commit path as the targets. Suyamoon's two
patches removed two fast-commit fallback reasons,
\texttt{XATTR}~(inline) and \texttt{FALLOC\_RANGE}, cutting JBD2
full commits by 64$\times$ and 62.5$\times$ on their respective
microbenches and reproducing the architectural claim
byte-identically on bare-metal. Sahil's CAS replacement of
\texttt{jbd2\_log\_wait\_commit} removed most of the
\texttt{j\_state\_lock} contention his \texttt{kretprobes}
measurement had identified, giving 40--54\% IOPS gains at 16
threads on three workloads. Shrey's FJM module replaced the
strict ring-buffer journal allocator with a free-block bitmap,
cutting P99 latency on the WAL workload by 8$\times$ in
\texttt{ordered} mode at the same WAF.

The shared methodological lesson is the same across all four
parts: \emph{measure the hit rate of the slow path you intend to
optimize before writing kernel code.} Suyamoon's C1 and C2 did
not do that and produced no measurable improvement. Milan's
characterization, Suyamoon's C3 and C4, Sahil's CAS module, and
Shrey's FJM all did, and all produce results that reproduce on
independent hardware.

% =====================================================================
\begin{thebibliography}{9}

\bibitem{ext4-design}
A. Mathur, M. Cao, S. Bhattacharya, A. Dilger, A. Tomas, L. Vivier.
\textit{The new ext4 filesystem: current status and future plans}.
Proceedings of the Linux Symposium, Ottawa, 2007.

\bibitem{ijournaling}
D. Park, D. Shin.
\textit{iJournaling: Fine-Grained Journaling for Improving the
Latency of Fsync System Call}. USENIX ATC 2017.

\bibitem{optfs}
V. Chidambaram, T. S. Pillai, A. C. Arpaci-Dusseau, R. H.
Arpaci-Dusseau. \textit{Optimistic Crash Consistency}. ACM SOSP 2013.

\bibitem{atc24-fc}
H. Shirwadkar, S. Kadekodi, T. Ts'o.
\textit{FastCommit: Resource-Efficient, Performant, and
Cost-Effective File System Journaling}. USENIX ATC 2024.

\bibitem{cjfs}
Y. Oh, J. Yoo, D. Nam, C. Min, Y. Won.
\textit{CJFS: Concurrent Journaling for Better Scalability}.
USENIX FAST 2023.

\bibitem{syncsync}
Q. Jiang et al.\
\textit{Sync+Sync: A Covert Channel Built on fsync with Storage}.
USENIX Security 2024.

\end{thebibliography}

% ========================================================
% APPENDIX
% ========================================================

\appendix
\clearpage
\section{Artifact Evaluation}

This appendix follows the CS614 Artifact Evaluation template.

\subsection{Artifact directory structure}

The submitted artifact is the ZIP \texttt{MTech Rocks.zip}, which
unpacks to a single folder \texttt{MTech Rocks/} with one
subdirectory per group member plus the umbrella documents at the
top:

\begin{lstlisting}[basicstyle=\ttfamily\scriptsize]
MTech Rocks/
|-- README.md, declaration.txt, project_report.pdf
|-- suyamoon/   (241110091): C1-C4 patches, benches, results, instructions
|-- sahil/      (241110061): jbd2_trace module, scaling matrix, results
|-- milan/      (241110042): WAL + combined fio harnesses, averaged tables
|-- shrey/      (251110068): FJM patches + ext4_tracker module + benches
+-- combined_report/  (LaTeX source of project_report.pdf)
\end{lstlisting}

The Linux kernel source tree is not in the ZIP. Reviewers should
fetch the pristine 6.1.4 tarball from
\url{https://cdn.kernel.org/pub/linux/kernel/v6.x/linux-6.1.4.tar.xz}
and apply the relevant patches to it.

\subsection{Setup}

\paragraph{Hardware.} 2~cores minimum (8+ recommended for Sahil's
matrix); 4~GiB RAM minimum (8~GiB recommended); 40~GB free disk
plus a dedicated ext4-formattable partition for Sahil
(default \texttt{/dev/nvme0n1p7}), Milan (\texttt{nvme0n1p6}),
and Shrey (\texttt{/dev/sdb}). No GPU.

\paragraph{Operating system.} Ubuntu 22.04 or 24.04 LTS. Tested on
Ubuntu 24.04.2 inside VirtualBox (Suyamoon VM data) and on a
bare-metal Ubuntu host with NVMe SSD (Suyamoon bare-metal data,
Sahil, Milan, Shrey).

\paragraph{Software dependencies.}
\begin{lstlisting}[language=bash]
sudo apt install -y fio attr build-essential libncurses-dev bison \
    flex libssl-dev libelf-dev bc dwarves zstd patch wget gawk \
    linux-tools-common linux-tools-$(uname -r) bpftrace sysbench \
    fsmark dbench python3-numpy python3-matplotlib blktrace \
    linux-headers-$(uname -r)
\end{lstlisting}

\paragraph{Kernel build.}
\begin{lstlisting}[language=bash]
wget https://cdn.kernel.org/pub/linux/kernel/v6.x/linux-6.1.4.tar.xz
tar xf linux-6.1.4.tar.xz && cd linux-6.1.4

# Suyamoon's two patches:
patch -p1 < /path/to/MTech\ Rocks/suyamoon/fc-inline-xattr.patch
patch -p1 < /path/to/MTech\ Rocks/suyamoon/fc-fallocate-range.patch

# Shrey's two patches + module (only needed if running C6):
git apply /path/to/MTech\ Rocks/shrey/patches/First.patch
git apply /path/to/MTech\ Rocks/shrey/patches/Second.patch
tar xzf /path/to/MTech\ Rocks/shrey/module/ext4_tracker.tar.gz -C fs/

cp /boot/config-$(uname -r) .config
sed -i 's/^CONFIG_LOCALVERSION=.*/CONFIG_LOCALVERSION="-mtechrocks"/' .config
make menuconfig    # for Shrey: File systems -> Ext4 TRACKER filesystem -> M
make olddefconfig
make -j$(nproc)
sudo make modules_install && sudo make install && sudo update-grub && sudo reboot
\end{lstlisting}

Sahil's module is built out-of-tree against any installed
6.1.4 with \texttt{linux-headers}; no in-tree patch required.
Milan's harness needs no custom kernel.

\subsection{Features supported}

\begin{itemize}[leftmargin=*]
  \item \textbf{Suyamoon:} C3 (inline xattr fast commit) and C4
        (fallocate range fast commit). Patches under
        \pth{suyamoon/}.
  \item \textbf{Sahil:} C5 (CAS-gated lockless
        \texttt{jbd2\_log\_wait\_commit}) plus a JBD2 lock-overhead
        tracer. Module under \pth{sahil/module/}.
  \item \textbf{Milan:} per-mode JBD2 characterization on WAL and
        WAL+sequential workloads. Harnesses under
        \pth{milan/benchmarks/}.
  \item \textbf{Shrey:} C6 (Fragmented Journal Map). Patches under
        \pth{shrey/patches/}; module source in
        \pth{shrey/module/ext4_tracker.tar.gz}.
\end{itemize}

\paragraph{Findings during evaluation.} No crashes, deadlocks, or
assertion failures were observed during any member's evaluation
of the patched kernel. \texttt{xfstests generic/473} fails
identically on stock 6.1.4 and on every patched kernel we tried;
this is a pre-existing failure, not a regression. Suyamoon's C3
and C4 patches compose cleanly. Shrey's FJM does not yet
implement the JBD2 recovery-path read of the on-disk index;
documented as a known limitation.

\subsection{Assumptions and unsupported}

C3 covers inline xattrs only (large block-stored xattrs and
\texttt{ea\_inode} xattrs still fall back). C4 covers
\texttt{COLLAPSE\_RANGE} and \texttt{INSERT\_RANGE} only;
\texttt{PUNCH\_HOLE} is unaffected. Sahil's module is not measured
against \texttt{data=journal} (SSD-bound). Shrey's FJM does not
yet consume the on-disk index on recovery.

\subsection{Getting started (within 30 minutes)}

\begin{enumerate}[leftmargin=*]
  \item Extract the ZIP. Verify Suyamoon's patches apply cleanly
        (\texttt{patch -p1 --dry-run} against a pristine 6.1.4
        tree).
  \item Build Sahil's module against your running kernel
        (\texttt{cd sahil/module \&\& make}).
  \item Inspect recorded numbers without running:
\begin{lstlisting}[language=bash]
cat "MTech Rocks/suyamoon/eval_results_c3/baremetal/PATCHED_6.1.4-C3Patch/xattr_loop.txt"
cat "MTech Rocks/suyamoon/eval_results_c4/baremetal/PATCHED_C4_6.1.4-C3C4Patch/collapse_summary.txt"
ls "MTech Rocks/sahil/results/benchmark_results/" | head
cat "MTech Rocks/milan/benchmarks/wal_workload_no_fc/sync_heavy_bs=4k.txt" | head -30
cat "MTech Rocks/shrey/results/logs/full_benchmark1.log" | head -20
\end{lstlisting}
  \item Each per-member \texttt{README.md} has the full reproduction
        recipe.
\end{enumerate}

\subsection{Detailed evaluation}

The table below lists every experiment in the artifact: purpose,
how to run, expected runtime, expected result, and where the output
lands. The table uses \texttt{longtable} so it breaks across pages
cleanly.

\begin{small}
\setlength{\LTpre}{0pt}
\setlength{\LTpost}{0pt}
\renewcommand{\arraystretch}{1.15}
\begin{longtable}{@{}p{0.4cm}p{4.0cm}p{3.7cm}p{7.5cm}@{}}
\toprule
\textbf{\#} & \textbf{Scenario} & \textbf{Script} & \textbf{Objective / Expected outcome} \\
\midrule
\endfirsthead

\multicolumn{4}{l}{\small\textit{(continued from previous page)}} \\
\toprule
\textbf{\#} & \textbf{Scenario} & \textbf{Script} & \textbf{Objective / Expected outcome} \\
\midrule
\endhead

\midrule
\multicolumn{4}{r}{\small\textit{(continued on next page)}} \\
\endfoot

\bottomrule
\endlastfoot

\multicolumn{4}{@{}l}{\textit{Suyamoon -- Candidate 3: inline xattr fast commit}} \\
1 & Primary xattr workload. 5000 \texttt{fsetxattr+fsync} ops; 256 MB loop fs; \texttt{-O fast\_commit}; single inline-sized value.
  & \pth{suyamoon/bench_xattr.sh}
  & Measure per-op latency and JBD2 commit count. Patched: $\sim$23~s, tx\_delta $\sim$78. Stock: $\sim$31~s, tx\_delta 5000. 64$\times$ commit reduction; 27\% wall-time reduction. \\

2 & Non-xattr control. fio 4-job randwrite+fsync for 30~s.
  & \pth{suyamoon/bench_async_prefetch.sh}
  & Confirm unrelated workloads are not regressed. Patched IOPS within $\pm$5\% of stock. Observed: 528 vs 530. \\

3 & Crash recovery A. 100 \texttt{setfattr} ops on one file; lazy umount; remount; count.
  & \pth{suyamoon/c3_crash_test_a.sh}
  & Expanded \texttt{FC\_TAG\_INODE} replay restores all inline xattrs. PASS with 100/100. \\

4 & Crash recovery B. Set 100; remove 50 even-numbered; simulated crash.
  & \pth{suyamoon/c3_crash_test_b.sh}
  & Removal path replays correctly via fast commit. PASS with exactly 50 odd xattrs remaining. \\

5 & Crash recovery C. 668-byte xattr (forces the block path).
  & \pth{suyamoon/c3_crash_test_c.sh}
  & Full-commit fallback is taken for block xattrs and data is preserved. PASS with all 668 bytes recovered. \\

\multicolumn{4}{@{}l}{\textit{Suyamoon -- Candidate 4: fallocate range fast commit}} \\
6 & Characterization on C3 kernel. 3 microbenches (fallocate, large-xattr block, concurrent fsync).
  & \pth{suyamoon/bench_fallocate_range.sh}, \pth{bench_xattr_block.sh}, \pth{bench_concurrent_fsync.sh}
  & Fallocate and xattr block show tx/op $\approx$ 1.001 (strong); concurrent fsync drops to 0.003 at T=16 (weak). Winner: fallocate. \\

7 & Primary fallocate workload. 1000 \texttt{fallocate+fsync} ops for each of COLLAPSE\_RANGE and INSERT\_RANGE.
  & \pth{suyamoon/bench_fallocate_range.sh}
  & C4 patched: tx\_delta 16, ms/op 5.2--5.4. C3 baseline (arch-equivalent to stock for this path): tx\_delta 1001, ms/op 7.0. 62.5$\times$ tx reduction, $\sim$24\% wall-time reduction. \\

8 & Crash recovery A (C4). 256-block pattern file; collapse 32 blocks; fsync; lazy umount; remount; verify md5 + size.
  & \pth{suyamoon/c4_crash_test_a.sh}
  & FC replay of ADD\_RANGE+DEL\_RANGE+FC\_TAG\_INODE reproduces expected post-collapse state. PASS. \\

9 & Crash recovery B (C4). Insert 16 blocks at offset 16; verify head + hole + shifted tail.
  & \pth{suyamoon/c4_crash_test_b.sh}
  & PASS (md5 + size match). \\

10 & Crash recovery C (C4). Three interleaved collapse/insert ops + crash.
  & \pth{suyamoon/c4_crash_test_c.sh}
  & PASS (md5 + size match the pure-Python expected state). \\

\multicolumn{4}{@{}l}{\textit{Suyamoon -- regression gate, applies to both patches}} \\
11 & xfstests subset \texttt{generic/\{062,118,300,337,454,455,473,482\}} + \texttt{ext4/032}.
  & xfstests
  & 6/6 runnable tests pass identically on stock, C3, and C4. One pre-existing \texttt{generic/473} failure on all three (not our regression). Logs at \pth{suyamoon/xfstests_c3/}, \pth{suyamoon/xfstests_c4/}. \\

\multicolumn{4}{@{}l}{\textit{Sahil -- Candidate 5: lockless CAS for jbd2\_log\_wait\_commit}} \\
12 & Lock-overhead measurement. Load module with \texttt{optimize=0}; run fio 16-thread randwrite+fsync; rmmod prints lock-time fraction to dmesg.
  & \pth{sahil/benchmarks/dmesg_sampler.sh}
  & $\sim$78\% of total fsync time spent in \texttt{jbd2\_log\_wait\_commit}. Recorded in \pth{sahil/results/dmesg_logs/}. \\

13 & Crash-safety smoke under the lockless module. Concurrent fsync flood + immediate umount + \texttt{fsck.ext4 -n}.
  & \pth{sahil/benchmarks/validation.sh}
  & fsck reports clean. \pth{sahil/results/fs_log.txt}. \\

14 & Scaling matrix. ordered/writeback $\times$ opt 0/1 $\times$ threads $\{1,4,8,16\}$ $\times$ 3 runs $\times$ \{fio, sysbench, fs\_mark, dbench\}.
  & \pth{sahil/benchmarks/run_evals.sh}
  & +40\% IOPS lock-contention, +54\% IOPS metadata, +7\% OLTP at 16 threads. Per-run JSON+txt under \pth{sahil/results/benchmark_results/}. \\

\multicolumn{4}{@{}l}{\textit{Milan -- workload characterization}} \\
15 & WAL workload across 4 modes. fio single-thread, \texttt{fdatasync}/write, block sizes $\{4,8,16,32,64\}$ KiB, 5 runs averaged.
  & \pth{milan/benchmarks/wal_workload_no_fc/run_repeated.sh}
  & Journaling halves IOPS regardless of mode; 99.6\% of latency is fdatasync. Averaged tables saved as \texttt{*.txt}. \\

16 & Concurrent WAL + sequential workload. Two fio jobs simultaneously bounded to the same wall time; bpftrace on JBD2 throughout.
  & \pth{milan/benchmarks/combined_workload_no_fc/run_repeated.sh}
  & Sequential bandwidth mode-insensitive; WAL tail latency degrades sharply under contention. Averaged tables saved as \texttt{*.txt}. \\

\multicolumn{4}{@{}l}{\textit{Shrey -- Candidate 6: Fragmented Journal Map}} \\
17 & FJM functional check + 6-mode benchmark. Loads \texttt{ext4\_tracker.ko}; fresh-formats /dev/sdb (64 MB journal); runs 100 MB WAL fio across ordered/journal/writeback $\times$ \{std,fjm\}.
  & \pth{shrey/benchmarks/fjm_full_benchmark.sh}
  & $\sim$40-60 min. \texttt{ordered\_fjm} P99 $\sim$0.05 ms vs $\sim$0.4 ms std at same WAF; \texttt{journal\_fjm} WAF $\sim$3.61$\times$ vs 3.50$\times$ std (index-block cost). \\

18 & FJM 5-run stability check.
  & \pth{shrey/benchmarks/run_repeated_benchmark.sh}
  & $\sim$3.5-5 hours. WAF stable across 5 runs in every mode. \\

\end{longtable}
\end{small}

\end{document}
